\section{Chapter 2 Background Study}

\subsection{2.1 Cloud Computing}

Cloud computing is a computing paradigm that enables users to access configurable computing resources---including servers, storage, databases, networking, software, analytics, and artificial intelligence services---over the internet without owning or maintaining physical infrastructure. These resources are provisioned on demand and can be rapidly scaled according to organizational requirements, allowing businesses to optimize operational costs while improving service availability and flexibility. Unlike traditional on-premises computing environments, cloud computing follows a utility-based service model in which organizations pay only for the resources they consume. This eliminates significant upfront investments in hardware procurement, infrastructure maintenance, and system administration. Cloud Service Providers (CSPs) such as Amazon Web Services (AWS), Microsoft Azure, and Google Cloud Platform (GCP) manage the underlying physical infrastructure while customers focus primarily on deploying and managing their applications. Cloud computing has become the foundation of modern digital transformation across various sectors, including finance, healthcare, education, manufacturing, government, and e-commerce. Organizations increasingly rely on cloud platforms because they offer high availability, disaster recovery capabilities, global scalability, automated resource provisioning, and seamless integration with advanced technologies such as artificial intelligence, machine learning, Internet of Things (IoT), and big data analytics. The National Institute of Standards and Technology (NIST) defines cloud computing as ``a model for enabling ubiquitous, convenient, on-demand network access to a shared pool of configurable computing resources that can be rapidly provisioned and released with minimal management effort or service provider interaction.'' This definition remains one of the most widely accepted foundations for academic research in cloud computing. Cloud computing is generally characterized by five essential characteristics:

\begin{itemize}
\tightlist
\item
  \textbf{On-demand Self-Service:} Users can independently provision computing resources without requiring manual intervention from the cloud provider.\\
\item
  \textbf{Broad Network Access:} Services are accessible over standard network protocols using various client devices such as laptops, smartphones, and workstations.\\
\item
  \textbf{Resource Pooling:} Physical and virtual resources are shared among multiple customers using a multi-tenant architecture while maintaining logical isolation.\\
\item
  \textbf{Rapid Elasticity:} Computing resources can be automatically scaled up or down based on changing workload demands.\\
\item
  \textbf{Measured Service:} Resource usage is continuously monitored, measured, and billed according to actual consumption.
\end{itemize}

These characteristics have significantly accelerated cloud adoption; however, they have also introduced new security challenges associated with dynamic resource provisioning, automated deployments, and increasingly complex cloud infrastructures.

\subsection{2.2 Cloud Service Models}

Cloud service providers deliver computing resources through multiple service models that differ in the level of control shared between the provider and the customer. Understanding these models is essential because security responsibilities vary depending on the selected service model.

\subsubsection{Infrastructure as a Service (IaaS)}

Infrastructure as a Service provides virtualized computing resources such as virtual machines, storage, networking components, and operating systems. The cloud provider manages the underlying physical infrastructure, while customers are responsible for configuring operating systems, applications, identity management, firewalls, and security policies. Examples include:

\begin{itemize}
\tightlist
\item
  Amazon EC2\\
\item
  Azure Virtual Machines\\
\item
  Google Compute Engine
\end{itemize}

IaaS offers maximum flexibility but also requires customers to manage a larger portion of the security stack.

\subsubsection{Platform as a Service (PaaS)}

Platform as a Service provides a managed application development environment where developers can build, test, and deploy applications without managing operating systems or hardware. Examples include:

\begin{itemize}
\tightlist
\item
  AWS Elastic Beanstalk\\
\item
  Azure App Service\\
\item
  Google App Engine
\end{itemize}

The cloud provider manages the infrastructure, runtime environment, middleware, and operating system, while customers focus primarily on application development and data management.

\subsubsection{Software as a Service (SaaS)}

Software as a Service delivers complete applications over the internet. Users simply access the software through a web browser while the cloud provider manages infrastructure, application updates, security patches, and maintenance. Examples include:

\begin{itemize}
\tightlist
\item
  Microsoft 365\\
\item
  Google Workspace\\
\item
  Salesforce
\end{itemize}

Although SaaS significantly reduces operational complexity, organizations remain responsible for managing user identities, access permissions, and sensitive organizational data.

\subsection{2.3 Cloud Deployment Models}

Cloud infrastructures can be deployed using different deployment models depending on organizational requirements for scalability, regulatory compliance, and security.

\subsubsection{Public Cloud}

In a public cloud environment, computing resources are owned and managed by third-party cloud providers and shared among multiple customers. Public clouds offer excellent scalability and cost efficiency, making it the preferred choice for startups and enterprises alike. Examples include AWS, Azure, and Google Cloud Platform.

\subsubsection{Private Cloud}

Private cloud infrastructure is dedicated exclusively to a single organization. It may be hosted on-premises or managed by a third-party provider. Private clouds provide greater control over security and regulatory compliance but typically require higher operational costs.

\subsubsection{Hybrid Cloud}

Hybrid cloud environments combine private and public cloud infrastructures, allowing organizations to keep sensitive workloads in private environments while utilizing public cloud services for scalability and cost optimization.

\subsubsection{Multi-Cloud}

A multi-cloud strategy involves using services from multiple cloud providers simultaneously. Organizations may combine AWS, Azure, and Google Cloud to improve resilience, avoid vendor lock-in, and optimize service availability. However, multi-cloud environments significantly increase management complexity and introduce additional security challenges due to inconsistent security policies across providers.

\subsection{2.4 Cloud Security Fundamentals}

Cloud security encompasses the technologies, policies, processes, and controls used to protect cloud-based infrastructure, applications, and data from unauthorized access, cyberattacks, accidental exposure, and operational failures. Unlike traditional IT environments, cloud security follows a Shared Responsibility Model, where security responsibilities are divided between the cloud provider and the customer. While the provider secures the underlying cloud infrastructure---including physical data centers, networking hardware, and virtualization platforms---the customer remains responsible for securing cloud resources such as identities, applications, operating systems, stored data, and access permissions. Core cloud security principles include:

\begin{itemize}
\tightlist
\item
  Identity and Access Management (IAM)\\
\item
  Network Security\\
\item
  Data Encryption\\
\item
  Security Monitoring and Logging\\
\item
  Vulnerability Management\\
\item
  Backup and Disaster Recovery\\
\item
  Compliance Management\\
\item
  Continuous Security Assessment
\end{itemize}

Failure to properly implement these security controls can lead to cloud misconfigurations that expose organizational assets to cyber threats.

\subsection{2.5 Identity and Access Management (IAM)}

Identity and Access Management (IAM) is one of the most critical security components within cloud environments. IAM controls who can access cloud resources, what actions they are permitted to perform, and under which conditions those actions are allowed. Modern cloud platforms implement IAM using several components, including users, groups, roles, policies, and permissions. Organizations assign permissions according to the principle of least privilege, ensuring that users receive only the minimum access necessary to perform their responsibilities. However, improperly configured IAM policies remain one of the leading causes \cite{vanede2022iam} of cloud security incidents. Examples include:

\begin{itemize}
\tightlist
\item
  Wildcard permissions (*)\\
\item
  Administrative access assigned to regular users\\
\item
  Long-lived access keys\\
\item
  Unused privileged accounts\\
\item
  Cross-account trust misconfigurations
\end{itemize}

Because IAM permissions determine access to virtually every cloud resource, attackers frequently target identity systems during cloud attacks. Consequently, CloudSentinel AI places significant emphasis on analyzing IAM relationships and identifying privilege escalation opportunities.

\subsection{2.6 Cloud Misconfigurations}

Cloud misconfiguration refers to the incorrect, incomplete, or insecure configuration of cloud resources, services, or security controls that unintentionally expose an organization's cloud infrastructure to security threats. Unlike traditional software vulnerabilities, which originate from programming defects, cloud misconfigurations primarily arise due to human error, improper deployment practices, inadequate security knowledge, or incorrect implementation of cloud services. Modern cloud platforms provide thousands of configurable parameters across networking, storage, identity management, databases, monitoring, encryption, and application deployment. While this flexibility enables organizations to build highly customized cloud infrastructures, it also increases the likelihood of accidental configuration mistakes. Even a single incorrectly configured cloud resource may allow attackers to gain unauthorized access, perform privilege escalation, exfiltrate sensitive information, or compromise multiple interconnected services.

Cloud misconfigurations have consistently been identified as one of the leading causes of cloud security incidents. Industry reports indicate that a significant proportion of successful cloud breaches result not from sophisticated zero-day exploits but from improperly configured cloud environments. As organizations increasingly adopt Infrastructure-as-Code (IaC), DevOps pipelines, serverless computing, and automated cloud provisioning, the speed of deployment often exceeds the speed of security validation, making configuration management a critical aspect of cloud security. Unlike software vulnerabilities that require attackers to exploit code-level weaknesses, misconfigured cloud resources are frequently exposed directly to the internet and may already provide legitimate access paths for malicious actors. Consequently, identifying and correcting cloud misconfigurations has become one of the primary objectives of modern cloud security programs.

\subsection{2.7 Common Cloud Misconfigurations in AWS}

Amazon Web Services provides hundreds of cloud services, each with numerous configurable security settings. Although AWS offers secure default configurations for many services, improper modifications or deployment errors can introduce significant security risks. Some of the most common AWS cloud misconfigurations include:

\begin{itemize}
\tightlist
\item
  \textbf{Publicly Accessible Amazon S3 Buckets:} Amazon S3 is widely used for storing application data, backups, multimedia files, and sensitive organizational documents. Improper bucket policies or Access Control Lists (ACLs) may unintentionally expose confidential data to the public internet. Numerous large-scale data breaches have occurred due to publicly accessible S3 buckets containing customer information, financial records, source code, or intellectual property.\\
\item
  \textbf{Overly Permissive IAM Policies:} Identity and Access Management (IAM) controls user permissions within AWS. Excessive permissions, wildcard actions (Action: ``*'') or unrestricted resource access (Resource: ``*'') violate the Principle of Least Privilege and significantly increase the attack surface. Attackers who compromise a low-privileged account may leverage excessive IAM permissions to escalate privileges and obtain administrative access.\\
\item
  \textbf{Misconfigured Security Groups:} AWS Security Groups function as virtual firewalls that regulate inbound and outbound network traffic for cloud resources. Allowing unrestricted inbound access (for example, opening SSH on port 22 or RDP on port 3389 to 0.0.0.0/0) exposes compute resources directly to the internet, making them vulnerable to brute-force attacks, credential theft, and unauthorized remote access.\\
\item
  \textbf{Disabled Logging and Monitoring:} CloudTrail, CloudWatch, and AWS Config provide essential logging and monitoring capabilities that enable organizations to detect suspicious activities, investigate security incidents, and maintain compliance. Disabling these services significantly reduces visibility into cloud operations, making incident detection and forensic investigations considerably more difficult.\\
\item
  \textbf{Unencrypted Storage Resources:} Cloud providers offer encryption mechanisms for services such as Amazon S3, Elastic Block Store (EBS), and Relational Database Service (RDS). Failure to enable encryption at rest or in transit increases the risk of unauthorized data disclosure in the event of unauthorized access or storage compromise.\\
\item
  \textbf{Exposed Secrets and Credentials:} Developers occasionally store API keys, database passwords, authentication tokens, or AWS access keys directly within application source code, configuration files, Git repositories, or environment variables. Exposure of these secrets allows attackers to authenticate directly against cloud services without exploiting software vulnerabilities.\\
\item
  \textbf{Public Database Instances:} Databases containing customer records or business-critical information should generally remain isolated within private networks. Accidentally assigning public IP addresses or allowing unrestricted inbound access exposes database servers to unauthorized users, significantly increasing the likelihood of data breaches.
\end{itemize}

Although each of these configuration errors may initially appear independent, attackers frequently combine multiple low-severity misconfigurations into complete attack paths capable of compromising entire cloud environments.

\subsection{2.8 Causes of Cloud Misconfigurations}

Cloud misconfigurations rarely result from a single technical failure. Instead, they emerge from a combination of organizational, operational, and human factors. The primary causes include:

\begin{itemize}
\tightlist
\item
  \textbf{Human Error:} Cloud administrators frequently configure large numbers of cloud resources under strict deployment timelines. Simple mistakes, such as assigning incorrect permissions or enabling public access, can unintentionally expose critical assets.\\
\item
  \textbf{Lack of Cloud Security Expertise:} Cloud platforms introduce new security models that differ substantially from traditional on-premises environments. Insufficient understanding of IAM policies, networking configurations, encryption mechanisms, or cloud-specific services often leads to insecure deployments.\\
\item
  \textbf{Infrastructure-as-Code (IaC) Errors:} Modern organizations increasingly automate cloud deployments using Infrastructure-as-Code technologies such as Terraform, AWS CloudFormation, and Pulumi. While automation improves consistency, configuration mistakes embedded within deployment templates may be replicated across hundreds of cloud resources.\\
\item
  \textbf{Rapid DevOps Deployment:} Continuous Integration and Continuous Deployment (CI/CD) pipelines prioritize rapid software delivery. Without automated security validation, insecure configurations may reach production environments before undergoing adequate security assessment.\\
\item
  \textbf{Configuration Drift:} Cloud environments continuously evolve as administrators modify security groups, IAM policies, storage permissions, or networking configurations. Over time, these incremental changes may diverge from the organization's original secure baseline, introducing previously unnoticed security weaknesses.\\
\item
  \textbf{Complexity of Modern Cloud Environments:} Large organizations often operate thousands of interconnected cloud resources across multiple AWS accounts, regions, and services. Understanding the security implications of each configuration change becomes increasingly difficult as infrastructure complexity grows.
\end{itemize}

\subsection{2.9 Shared Responsibility Model}

One of the fundamental principles of cloud security is the Shared Responsibility Model, which defines the division of security responsibilities between the cloud service provider and the customer. Under this model, the cloud provider is responsible for securing the underlying cloud infrastructure, including physical data centers, networking equipment, virtualization layers, and managed cloud services. Customers, however, remain responsible for securing their workloads, identities, operating systems, applications, configurations, and organizational data.

For example, Amazon Web Services secures the physical infrastructure supporting services such as Amazon EC2 and Amazon S3. However, customers are responsible for configuring IAM policies, defining Security Group rules, enabling encryption, securing operating systems, patching applications, and protecting sensitive data stored within their cloud environments. This distinction is particularly important because many organizations incorrectly assume that cloud providers automatically secure all aspects of their cloud deployments. In reality, most cloud security incidents originate from customer-side configuration errors rather than vulnerabilities within the cloud provider's infrastructure. Understanding the Shared Responsibility Model is therefore essential for designing effective cloud security strategies and forms one of the primary motivations for automated cloud security posture management solutions.

\subsection{2.10 Cloud Security Posture Management (CSPM)}

Cloud Security Posture Management (CSPM) refers to a category of security solutions designed to continuously monitor cloud environments, identify security misconfigurations, evaluate compliance with security standards, and assist organizations in maintaining secure cloud deployments. CSPM platforms automatically collect configuration data from cloud providers through APIs and compare these configurations against predefined security policies, compliance frameworks, and industry best practices. Examples of commonly supported frameworks include the Center for Internet Security (CIS) Benchmarks, the National Institute of Standards and Technology (NIST) Cybersecurity Framework, ISO/IEC 27001, PCI DSS, HIPAA, and SOC 2.

Typical CSPM capabilities include:

\begin{itemize}
\tightlist
\item
  Continuous cloud asset discovery\\
\item
  Misconfiguration detection\\
\item
  Compliance assessment\\
\item
  Security policy validation\\
\item
  Risk reporting\\
\item
  Alert generation\\
\item
  Remediation recommendations
\end{itemize}

Several commercial and open-source CSPM solutions are widely adopted within industry, including AWS Security Hub, AWS Config, Microsoft Defender for Cloud, Wiz, Prisma Cloud, Prowler, ScoutSuite, and Steampipe. Although CSPM significantly improves cloud security visibility, most platforms primarily rely on rule-based detection mechanisms. Consequently, they often identify individual security findings without considering relationships among cloud resources or evaluating how multiple low-severity issues may collectively enable complex attack scenarios. This limitation represents one of the primary research motivations behind the proposed CloudSentinel AI framework.

\subsection{2.11 Cloud-Native Application Protection Platform (CNAPP)}

As cloud computing environments have evolved, traditional Cloud Security Posture Management (CSPM) solutions have expanded into a broader security architecture known as the Cloud-Native Application Protection Platform (CNAPP). CNAPP integrates multiple cloud security technologies into a unified platform capable of protecting cloud infrastructure, applications, workloads, identities, and data throughout the software development lifecycle. Unlike conventional CSPM solutions that primarily focus on identifying configuration errors, CNAPP combines Cloud Workload Protection Platforms (CWPP), Cloud Infrastructure Entitlement Management (CIEM), Kubernetes security, Infrastructure-as-Code (IaC) scanning, vulnerability management, runtime protection, and threat detection into a single security framework. This integrated approach provides organizations with comprehensive visibility across modern cloud-native environments. Despite these advancements, current CNAPP platforms continue to face several challenges. Many solutions generate an overwhelming number of security findings without effectively prioritizing risks based on organizational context. Furthermore, although they consolidate multiple security functions, most platforms still rely heavily on predefined detection rules and limited contextual reasoning. Consequently, security analysts must manually investigate how multiple alerts relate to one another before determining the overall security impact. The proposed CloudSentinel AI framework complements the objectives of CNAPP by introducing context-aware analysis, attack path modeling, and AI-assisted explanations that help security teams better understand and prioritize cloud security risks.

\subsection{2.12 Attack Graphs}

An attack graph is a graphical representation \cite{phillips1998graph,sheyner2002automated,jha2002formal} of the possible paths an attacker may exploit to compromise systems within a network or cloud environment. Rather than analyzing individual vulnerabilities independently, attack graphs model the relationships between assets, identities, permissions, network connectivity, and security controls to illustrate how an attacker could progress from an initial point of access toward critical organizational resources. In cloud environments, attack graphs have become increasingly important because cloud resources are highly interconnected. A seemingly minor configuration error may not represent a serious threat when viewed independently; however, when combined with excessive IAM permissions, exposed storage resources, weak network segmentation, or compromised credentials, it may enable complete attack chains capable of compromising sensitive cloud assets. For example, an attacker may initially exploit an internet-facing virtual machine exposed through an overly permissive security group. After gaining access, the attacker could abuse an overly permissive IAM role to retrieve credentials from AWS Secrets Manager, access confidential Amazon S3 buckets, and ultimately compromise sensitive organizational data. Although each individual configuration issue might be classified as medium severity, their combined exploitation path represents a critical organizational risk. Attack graphs therefore enable security analysts to understand not only what security issues exist but also how those issues interact to facilitate real-world attacks. This capability forms one of the fundamental design principles of the proposed CloudSentinel AI framework.

\subsection{2.13 Knowledge Graphs in Cloud Security}

Knowledge graphs provide a structured representation \cite{hogan2021knowledge} of entities and the relationships that exist between them. In cloud security, entities may include cloud services, virtual machines, IAM users, IAM roles, storage buckets, databases, security groups, network components, secrets, and organizational policies. Relationships describe how these entities interact, communicate, or depend upon one another. Unlike traditional relational databases that primarily store isolated records, knowledge graphs explicitly model interconnected relationships, making them highly suitable for representing complex cloud infrastructures. By capturing resource dependencies and security relationships, knowledge graphs enable more intelligent analysis of cloud environments. Within the proposed CloudSentinel AI framework, the knowledge graph serves as the central representation of the cloud environment. Cloud assets collected from AWS APIs are transformed into graph nodes, while permissions, network connections, ownership relationships, and resource dependencies are represented as graph edges. This graph-based representation allows the framework to identify privilege escalation paths, lateral movement opportunities, exposed assets, and multi-stage attack scenarios that may not be apparent through rule-based analysis alone. Furthermore, knowledge graphs support advanced graph algorithms capable of discovering hidden relationships among cloud resources, thereby improving contextual risk assessment and attack path generation.

\subsection{2.14 Context-Aware Risk Assessment}

Traditional cloud security tools generally assign predefined severity levels to individual security findings. For example, a publicly accessible storage bucket may always be classified as ``High Risk'' regardless of the type of data stored within the bucket or its role within the organization's infrastructure. However, real-world security risks are highly dependent upon operational context. Two seemingly identical cloud resources may present vastly different security implications depending on factors such as data sensitivity, network exposure, identity permissions, workload criticality, regulatory requirements, and business importance. Context-aware risk assessment extends conventional risk evaluation by incorporating these additional contextual factors into the overall security analysis. Instead of evaluating cloud resources independently, contextual analysis considers the relationships among multiple cloud assets and determines how combined security weaknesses influence the organization's overall attack surface. Within CloudSentinel AI, contextual risk assessment combines information obtained from configuration analysis, attack graph modeling, IAM relationships, network connectivity, asset criticality, and resource dependencies to generate dynamic risk scores that more accurately reflect the actual likelihood and impact of successful cyberattacks. This approach enables security teams to prioritize remediation efforts based on exploitability rather than relying solely on predefined severity classifications.

\subsection{2.15 Explainable Artificial Intelligence (XAI)}

Artificial Intelligence has become an increasingly important component of cybersecurity applications, including anomaly detection, malware classification, intrusion detection, and cloud security analysis. However, many AI models operate as ``black boxes,'' producing predictions without providing sufficient explanations regarding the reasoning behind their decisions. Explainable Artificial Intelligence (XAI) addresses this limitation \cite{rjoub2023xai,charmet2022xai} by developing AI systems capable of producing transparent, interpretable, and human-understandable explanations for their outputs. Rather than simply identifying a security issue, XAI explains why the issue is considered important, how it affects the organization, and which remediation actions should be prioritized. In cloud security, explainability is particularly valuable because security analysts must justify remediation decisions, communicate risks to management, and maintain compliance with organizational governance requirements. AI-generated explanations therefore improve both analyst productivity and organizational trust in automated security systems. Within CloudSentinel AI, the Explainable AI module will generate natural language descriptions for identified attack paths, contextual risk scores, potential attack scenarios, and recommended mitigation strategies. Instead of presenting only technical alerts, the framework will assist analysts in understanding the broader security implications of identified cloud misconfigurations.

\subsection{2.16 Large Language Models (LLMs) in Cloud Security}

Large Language Models (LLMs) have recently demonstrated \cite{vaswani2017attention,brown2020language,openai2023gpt4} significant potential across multiple cybersecurity domains, including vulnerability analysis, malware explanation, security log interpretation, incident response, and security policy generation. Their ability to understand natural language and synthesize technical information makes them valuable assistants for cloud security operations. Within cloud environments, LLMs can analyze security findings, summarize complex IAM policies, explain configuration errors, generate remediation recommendations, and assist analysts in interpreting cloud security alerts. However, existing implementations generally use LLMs only as auxiliary explanation tools rather than integrating them into broader contextual security analysis frameworks. The proposed CloudSentinel AI framework leverages LLM capabilities not for replacing existing CSPM solutions but for enhancing analyst understanding. After contextual risk assessment and attack path generation are completed, the LLM module produces concise, human-readable explanations describing the identified risks, potential attack progression, affected cloud resources, and recommended corrective actions. This approach reduces the cognitive burden on security analysts while improving the usability of cloud security findings.

\subsection{2.17 Need for CloudSentinel AI}

The increasing adoption of cloud computing has fundamentally transformed organizational infrastructure while simultaneously introducing new categories of security risks. Existing Cloud Security Posture Management and CNAPP solutions have significantly improved cloud visibility by automating the detection of configuration errors and compliance violations. Nevertheless, current platforms primarily focus on identifying isolated security findings rather than understanding how multiple cloud resources interact within complex cloud ecosystems. Modern cyberattacks rarely exploit a single misconfiguration. Instead, attackers combine multiple weaknesses---including identity misconfigurations, network exposure, insecure storage permissions, and excessive privileges---to establish attack paths that ultimately compromise sensitive organizational assets. Consequently, organizations require security solutions capable of reasoning about relationships among cloud resources instead of evaluating each configuration independently. CloudSentinel AI addresses this challenge by integrating cloud asset discovery, graph-based resource modeling, attack path analysis, contextual risk assessment, and Explainable Artificial Intelligence into a unified cloud security framework. Rather than producing extensive lists of disconnected alerts, the proposed system identifies the most critical attack paths, explains their potential business impact, and recommends prioritized remediation actions. By combining graph intelligence with AI-assisted security analysis, CloudSentinel AI aims to bridge the gap between conventional rule-based cloud security assessment and intelligent, context-aware cloud risk management. The framework is designed to improve security visibility, reduce alert fatigue, support faster incident response, and assist organizations in proactively securing increasingly complex cloud infrastructures.
